\documentclass[10pt,a4paper]{amsart}
\usepackage[margin=19mm]{geometry}
\usepackage{amsmath,amssymb,mathtools}
\usepackage[scr=rsfs,cal=euler]{mathalfa}
\usepackage{xfrac}
\usepackage[unicode,hidelinks,bookmarks=false]{hyperref}
\newcommand{\ie}{i.e.\ }
\newcommand{\cf}{cf.\ }
\newcommand{\resp}{respectively\ }
\newcommand{\Subsection}{\subsection}
\providecommand{\nobreakdash}{\nobreak\hbox{-}}
\setlength{\emergencystretch}{2em}
\multlinegap0pt
\newtheorem{proposition}{Proposition}
\title{Positivity of Second-Order Nearly Linearly Recurrent Sequences}
\author{Mahsa Shirmohammadi, James Worrell}
\date{Source: arXiv:2508.00944v4 (CC BY 4.0)}
\begin{document}
\maketitle

\noindent\emph{Source and licence.} Positivity of Second-Order Nearly Linearly Recurrent Sequences. Version arXiv:2508.00944v4 (CC BY 4.0), distributed by arXiv under the Creative Commons Attribution 4.0 International licence, http://creativecommons.org/licenses/by/4.0/, as recorded in the arXiv licence metadata for this exact version and shown on the abstract page https://arxiv.org/abs/2508.00944v4. Full licence notice: \texttt{licenses/arxiv-2508.00944-CC-BY-4.0.txt}.\par\smallskip

\noindent\textit{Editorial scope.} Counted unit: the article's proposition prop:non-vanish3 (source0.tex lines 3213--3244). The article's complete original proof of it is delivered with it (source0.tex lines 3246--3340, one \texttt{proof} environment of 95 lines). No mathematics is written, repaired or re-derived: the statement and the proof are the article's own lines, in the article's own notation, and no retained line is substituted or re-flowed.  Retained uncounted as the source-local context the proof needs: lines 3143--3147.  Nothing was omitted from any retained span, so the package carries no omission record.  No retained line contains a citation, so the wrapper carries no bibliography and no bibliography entry.  No retained line contains a figure environment, an image-inclusion command or drawing code, so no image is delivered and the package declares no asset dependency.  Editorial formatting only, in addition: an AMS \texttt{amsart} preamble that restates the article's own declarations and macros used by the retained lines, namely the environments \texttt{proposition} on separate counters, together with the standard packages those lines need.  The retained environments print in this document as Proposition 1 in order of appearance, whereas the article prints the same items with the numbering of its own full document.  The complete native source of the article as distributed in the arXiv e-print for this exact version is bundled unchanged as \texttt{sources/182584/source0.tex}.
\begin{equation}
    u_m
    = b(m)\lambda^m + \overline{b(m)\lambda^m} \, ,
    \label{eq:closed-form-repeated}
\end{equation}

\begin{proposition}
\label{prop:non-vanish3}
Fix $k\in\{0,\ldots,2d-1\}$. For $m,n\in\mathbb{N}$, let
\begin{equation}
%\begin{split}
    v_{n,m,k}
    :={a}_{n,2d-k}u_{m+kq_n}
       +\cdots
       +{a}_{n,2d-1}u_{m+q_n}
       +{a}_{n,2d}u_m
   % &=\sum_{h=0}^{k}a_{n,2d-h}u_{m+hq_n}.
%\end{split}
    \label{eq:def-vnmk}
\end{equation}
Then
\begin{equation}
    v_{n,m,k}
    =c_k(m,q_n)\lambda^m
     +\overline{c_k(m,q_n)\lambda^m},
    \label{eq:v-c-form}
\end{equation}
where $c_k(X,Y)$ is a non-zero exponential polynomial of the form
\begin{equation}
    c_k(X,Y)
    =\sum_{\ell=1}^{L_k}\sum_{i,j}
       \alpha_{k,\ell,i,j}X^iY^j\gamma_{k,\ell}^{\,Y},
    \label{eq:ck-exp-polynomial}
\end{equation}
where 
$\alpha_{k,\ell,i,j} \in \overline{\mathbb{Q}}$, and each
$\gamma_{k,\ell}$ is a monomial in $\lambda$ and $\overline{\lambda}$.
\end{proposition}

\begin{proof}
For $i\in\{0,\ldots,2d\}$ and $Y\in\mathbb{N}_0$, define
\begin{equation}
    A_i(Y)
    :=(-1)^i
    s_i\bigl(
        \underbrace{\lambda^Y,\ldots,\lambda^Y}_{d\text{ copies}},
        \underbrace{\overline{\lambda}^{\,Y},\ldots,
                    \overline{\lambda}^{\,Y}}_{d\text{ copies}}
    \bigr).
    \label{eq:def-Ai}
\end{equation}
Then $A_i(q_n)=a_{n,i}$. Moreover, the multiset of variables in
\eqref{eq:def-Ai} is invariant under complex conjugation, so
 $a_{n,i}$ is real for every $n$ and $i$.

Define
\begin{equation}
    c_k(X,Y)
    :=\sum_{h=0}^{k}
       A_{2d-h}(Y)b(X+hY)\lambda^{hY}.
    \label{eq:def-ck}
\end{equation}
Using \eqref{eq:closed-form-repeated} and
$A_i(q_n)=a_{n,i}$, we obtain
\begin{align*}
    v_{n,m,k}
    &=\sum_{h=0}^{k}A_{2d-h}(q_n)u_{m+hq_n}\\
    &=\sum_{h=0}^{k}A_{2d-h}(q_n)
      \left(
          b(m+hq_n)\lambda^{m+hq_n}
          +\overline{b(m+hq_n)\lambda^{m+hq_n}}
      \right)\\
    &=\lambda^m
      \sum_{h=0}^{k}A_{2d-h}(q_n)
          b(m+hq_n)\lambda^{hq_n}
      +\overline{
          \lambda^m
          \sum_{h=0}^{k}A_{2d-h}(q_n)
              b(m+hq_n)\lambda^{hq_n}
       }\\
    &=c_k(m,q_n)\lambda^m
      +\overline{c_k(m,q_n)\lambda^m}.
\end{align*}
This proves \eqref{eq:v-c-form}.

It remains to establish the asserted form and the non-vanishing of $c_k$.
For every $i\in\{0,\ldots,2d\}$, choosing $r$ of the $d$ copies of
$\lambda^Y$ and $s$ of the $d$ copies of
$\overline{\lambda}^{\,Y}$ gives
\begin{equation}
    A_i(Y)
    =(-1)^i
      \sum_{\substack{r+s=i\\0\leq r,s\leq d}}
      \binom{d}{r}\binom{d}{s}
      \bigl(\lambda^r\overline{\lambda}^{\,s}\bigr)^Y.
    \label{eq:Ai-expansion}
\end{equation}
Substituting \eqref{eq:Ai-expansion} into \eqref{eq:def-ck} yields
\begin{equation}
\begin{split}
    c_k(X,Y)
    ={}&\sum_{h=0}^{k}(-1)^{2d-h}
       \sum_{\substack{r+s=2d-h\\0\leq r,s\leq d}}
       \binom{d}{r}\binom{d}{s}
       b(X+hY)
       \bigl(\lambda^{h+r}\overline{\lambda}^{\,s}\bigr)^Y.
\end{split}
    \label{eq:ck-expanded}
\end{equation}
Since $b(X+hY)\in\overline{\mathbb{Q}}[X,Y]$, expression
\eqref{eq:ck-expanded} has the form \eqref{eq:ck-exp-polynomial}, and each
of its exponential bases is a monomial in $\lambda$ and
$\overline{\lambda}$.

Finally, evaluate \eqref{eq:def-ck} at $Y=0$. All the variables occurring
in \eqref{eq:def-Ai} are then equal to $1$, and hence
\(A_i(0)=(-1)^i\binom{2d}{i}\).
Therefore
\begin{align*}
    c_k(X,0)
    &=b(X)\sum_{h=0}^{k}A_{2d-h}(0)\\
    &=b(X)\sum_{h=0}^{k}(-1)^h\binom{2d}{h}\\
    &=(-1)^k\binom{2d-1}{k}b(X)\, ,
\end{align*}
where the last equality is the standard partial alternating-binomial
identity
\[
    \sum_{h=0}^{k}(-1)^h\binom{N}{h}
    =(-1)^k\binom{N-1}{k},
    \qquad 0\leq k<N.
\]
Since $k\leq 2d-1$ and $b$ is nonzero, $c_k(X,0)$ is nonzero. Thus
$c_k(X,Y)$ is not the zero exponential polynomial.
\end{proof}

\end{document}
